% Generated draft: requires mathematical and source review.
\documentclass[11pt]{article}
\usepackage[margin=25mm]{geometry}
\usepackage{fontspec}
\setmainfont{Latin Modern Roman}
\usepackage{amsmath,amssymb,mathtools,unicode-math}
\setmathfont{Latin Modern Math}
\usepackage{enumitem,needspace,xurl,hyperref}
\setlength{\parindent}{0pt}
\setlength{\parskip}{5pt}
\setlength{\emergencystretch}{3em}
\newcommand{\N}{\mathbb{N}}
\newcommand{\Z}{\mathbb{Z}}
\newcommand{\Q}{\mathbb{Q}}
\newcommand{\R}{\mathbb{R}}
\newcommand{\C}{\mathbb{C}}
\newcommand{\F}{\mathbb{F}}
\newcommand{\A}{\mathbb{A}}
\newcommand{\PP}{\mathbb{P}}
\newcommand{\E}{\mathbb{E}}
\newcommand{\eps}{\varepsilon}
\newcommand{\dd}{\,\mathrm{d}}
\begin{document}
\section*{UnsolvedMath ID: 20003112}
\textbf{Source problem code:} catalogue entry
\subsection{Definitions and mathematical statement}
The actual question attributed to D. Dumas is whether one can define a Thurston metric on universal Teichm\"uller space. The following section heading about infinitely generated first-kind Fuchsian groups and its introductory sentence are record-boundary contamination, not hypotheses of Problem 3.9. Use the model
\[
\mathcal T_{\rm univ}=\operatorname{Mob}^+(S^1)\backslash\operatorname{QS}^+(S^1),
\]
where equivalence is postcomposition by a M\"obius transformation. The source does not state which axioms, finiteness, topology, geodesics or compatibility with finite-type Thurston metrics are required.

The supplied aid proposes a metric on a proper sublocus. Let \(\mathcal G\) be the space of unoriented geodesics of \(\mathbb H^2\), identified with distinct unordered boundary pairs, and let \(L_0\) be its standard Liouville current. For a quasisymmetric boundary map \(f\), let \(L_f=f^*L_0\); this is unchanged under postcomposition by M\"obius maps. Define
\[
\mathcal U=\{[f]:\ \exists\,0<c_f\le C_f<\infty,\quad c_fL_0\le L_f\le C_fL_0\}.
\]
The inequalities are domination of measures, not bounds on finitely many selected boxes. On this locus set
\[
d_{\rm LD}([f],[g])=\log\operatorname*{ess\,sup}_{L_0}
\frac{\dd L_g/\dd L_0}{\dd L_f/\dd L_0}.
\]
The aid reports that this is a finite directed metric: it is nonnegative, separates equivalence classes, and satisfies the triangle inequality. For the stated Liouville normalization, complementary boxes \(Q,Q^\perp\) determined by four cyclic boundary points obey
\[
e^{-L_f(Q)}+e^{-L_f(Q^\perp)}=1.
\]
This rules out strict one-sided domination with constant below one and forces equality of currents in the zero-distance case; equality of pullback currents determines the boundary map up to M\"obius postcomposition. The Radon--Nikodym quotient chain rule supplies the triangle inequality. The report stresses that \(\mathcal U\) is proper: quasisymmetric power maps \(x\mapsto\operatorname{sgn}(x)|x|^a\) on the real boundary, \(a>0\), \(a\ne1\), supply failures of uniform comparison. Thus the construction is not a finite metric on all of universal Teichm\"uller space.

Two literal transplants are also obstructed. The disk has no essential closed-curve length family. Moreover unrestricted homeomorphic boundary extensions do not give a useful forward Lipschitz constant. In hyperbolic polar coordinates, for \(0<a\le1\),
\[
R_a(r,\theta)=(ar,\theta)
\]
is a disk homeomorphism extending the identity boundary map, with hyperbolic Lipschitz constant at most \(a\). Hence the infimum is zero already for the identity boundary datum; composing with a finite-Lipschitz reference extension gives the same collapse whenever such an extension is available. Taking its logarithm cannot produce a normalized distance. These unrestricted maps need not belong to a more tightly controlled extension class. The aid's proper-locus metric and these failures leave the general definition and desired compatibility properties unresolved.
\subsection{Short English statement}
Seek an appropriate universal Thurston metric after removing the next section's text from the record. The supplied aid defines a finite Liouville-domination metric only on uniformly comparable currents and explains why empty closed-curve data and unrestricted forward-Lipschitz extensions do not supply a global answer.
\subsection{Sources}
\begin{itemize}
\item D. Dumas, Problem 3.9, AIM workshop \emph{Lipschitz metric on Teichmueller space}, \url{https://aimath.org/pastworkshops/lipschitzteichproblems.pdf}. The supplied verification identifies the exact boundary on PDF page 4 and excludes the subsequent first-kind Fuchsian-group section.
\item Local UnsolvedMath registry, permanent ID 20003112, exact code \texttt{catalogue entry}; canonical location \texttt{aim-topology-notes.json}, notes[199], dated 14 August 2026; statement labeled exact although its stored text includes the explicitly separated contamination.
\item Research attempt 1: partial result, problem status unknown at the run, rigor label proved, novelty confidence low; registry status partially solved. The comparable-current construction and radial/power-map tests are machine-generated research aids, not peer reviewed; independent verification is required before citation. No separate authoritative bibliography is supplied.
\end{itemize}
\end{document}
